\documentclass{article}
\usepackage{graphicx, ctex, amsmath, geometry, amssymb}

\title{\bfseries{\huge{数学作业}}}
\author{西安科技大学\hspace{1em}通信学院\hspace{1em}电子科学与技术 \\
学号：25207035001\hspace{1em}姓名：张凯文}
\date{2026年5月}
\linespread{1.5}
\geometry{a4paper, margin=2cm}

\begin{document}

\maketitle

\newpage
\section{第一题}\noindent
在矩阵空间 $\mathbb{R}^{2\times2}$ 中，给定矩阵 $B = \begin{bmatrix} 
1 & 2 \\
0 & 3
\end{bmatrix}$，线性变换 $T$ 为
\[
T(X)=BX-XB,\quad(\forall X \in \mathbb{R}^{2\times2})
\]
求 $N(T)$ 和 $R(T)$ 的基与维数。

\textbf{解：}
\begin{align*}
N(T) &= \left\{ X \in \mathbb{R}^{2\times2} \mid T(X)=0 \right\} \\
\text{设 } X &= \begin{bmatrix} a & b \\ c & d \end{bmatrix} \\
T(X) &= BX - XB = 
\begin{bmatrix} 1 & 2 \\ 0 & 3 \end{bmatrix}
\begin{bmatrix} a & b \\ c & d \end{bmatrix} -
\begin{bmatrix} a & b \\ c & d \end{bmatrix}
\begin{bmatrix} 1 & 2 \\ 0 & 3 \end{bmatrix} \\
&= \begin{bmatrix} c-3b & -a-2b+d \\ 3a+2c-3d & 3b-c \end{bmatrix} \\
\text{令 } T(X)=0 \Rightarrow &
\begin{cases}
c-3b = 0 \\
-a-2b+d = 0 \\
3a+2c-3d = 0 \\
3b-c = 0
\end{cases}
\Rightarrow
\begin{cases}
c = 3b \\
d = a+2b
\end{cases} \\
\Rightarrow X &= \begin{bmatrix} a & b \\ 3b & a+2b \end{bmatrix}
= a\begin{bmatrix} 1 & 0 \\ 0 & 1 \end{bmatrix} + b\begin{bmatrix} 0 & 1 \\ 3 & 2 \end{bmatrix} \\
\because &
\begin{pmatrix} 1 \\ 0 \\ 0 \\ 1 \end{pmatrix}
\text{ 与 }
\begin{pmatrix} 0 \\ 1 \\ 3 \\ 2 \end{pmatrix}
\text{ 线性无关}
\Rightarrow M_1 = \begin{bmatrix} 1 & 0 \\ 0 & 1 \end{bmatrix},\;
M_2 = \begin{bmatrix} 0 & 1 \\ 3 & 2 \end{bmatrix}
\text{ 线性无关} \\
\therefore & M_1, M_2 \text{ 是 } N(T) \text{ 的一组基，} \dim N(T) = 2 \\
R(T) &= \left\{ T(X) \mid X \in \mathbb{R}^{2\times2} \right\}
\end{align*}
\[
\because \dim N(T) + \dim R(T) = \dim \mathbb{R}^{2\times2} = 4
\quad\Rightarrow\quad \dim R(T) = 2
\]

\[
T(E_{11}) = \begin{bmatrix}0 & -2 \\ 0 & 0 \end{bmatrix},\quad
T(E_{12}) = \begin{bmatrix}0 & -2 \\ 0 & 0 \end{bmatrix},\quad
T(E_{21}) = \begin{bmatrix}2 & 0 \\ 2 & -2 \end{bmatrix},\quad
T(E_{22}) = \begin{bmatrix}0 & 2 \\ 0 & 0 \end{bmatrix}
\]
\[
\gamma_1=\begin{pmatrix} 0 \\ -2 \\ 0 \\ 0\end{pmatrix},\;
\gamma_2=\begin{pmatrix} 0 \\ -2 \\ 0 \\ 0\end{pmatrix},\;
\gamma_3=\begin{pmatrix} 2 \\ 0 \\ 2 \\ -2\end{pmatrix},\;
\gamma_4=\begin{pmatrix} 0 \\ 2 \\ 0 \\ 0\end{pmatrix}
\]
其中 $\gamma_1$ 与 $\gamma_3$ 线性无关，故 $R(T)$ 的一组基可取为 $T(E_{11})$ 和 $T(E_{21})$。

\newpage

\section{第二题}\noindent
设 $B = \begin{bmatrix} 1 & 1 \\ 0 & 1 \end{bmatrix}$，线性空间
\[
V = \left\{ X=(x_{ij})_{2\times2} \mid x_{11}+x_{22} = 0,\; x_{ij}\in \mathbb{R} \right\}
\]
中的线性变换为 $T(X)=B^{\mathsf{T}}X - X^{\mathsf{T}}B$（$\forall X \in V$），求 $T$ 的特征值和特征向量。

\textbf{解：}
\[
X = \begin{bmatrix} x_{11} & x_{12} \\ x_{21} & -x_{11} \end{bmatrix} = 
x_{11}\begin{bmatrix}1 & 0 \\ 0 & -1\end{bmatrix} + 
x_{12}\begin{bmatrix}0 & 1 \\ 0 & 0\end{bmatrix} + 
x_{21}\begin{bmatrix}0 & 0 \\ 1 & 0\end{bmatrix}\] \\
\[
\text{选一组基}X_1 = \begin{bmatrix}1 & 0 \\ 0 & -1\end{bmatrix}
X_2 = \begin{bmatrix}0 & 1 \\ 0 & 0\end{bmatrix}
X_3 = \begin{bmatrix}0 & 0 \\ 1 & 0\end{bmatrix}
\Rightarrow X = x_{11}X_1 + x_{12}X_2 + x_{21}X_3
\]
\[
T(X)=\begin{bmatrix}1 & 0 \\ 1 & 1\end{bmatrix}
\begin{bmatrix}x_{11} & x_{12} \\ x_{21} & x_{22}\end{bmatrix} - 
\begin{bmatrix}x_{11} & x_{12} \\ x_{21} & x_{22}\end{bmatrix}
\begin{bmatrix}1 & 1 \\ 0 & 1\end{bmatrix} =
\begin{bmatrix}0 & x_{12}-x_{11}-x_{21} \\ x_{11}+x_{21}-x_{12} & 0\end{bmatrix}
\]
\[
T(X_1) = \begin{bmatrix}0 & -1 \\ 1 &0\end{bmatrix}
T(X_2) = \begin{bmatrix}0 & 1 \\ -1 &0\end{bmatrix}
T(X_3) = \begin{bmatrix}0 & -1 \\ 1 &0\end{bmatrix}
\]
\[
T(X)=AX = A\begin{bmatrix}X_1 & X_2 & X_3\end{bmatrix}\]
\[T(X_1) = 0 · X_1 + (-1) · X_2 + 1 · X_3,
T(X_2) = 0 · X_1 + 1 · X_2 + (-1) · X_3,
T(X_3) = 0 · X_1 + (-1) · X_2 + 1 · X_3 \]
\[\Rightarrow T(X) = \begin{bmatrix}0 & 0 & 0 \\ -1 & 1 & -1 \\ 1 & -1 & -1\end{bmatrix}
\begin{bmatrix}X_1 & X_2 & X_3\end{bmatrix}\]

\[\mid \lambda E - A \mid \Rightarrow \lambda_1,\lambda_2 = 0, \lambda_3 = 2,
\alpha_1 = \begin{bmatrix}1 \\ 1 \\ 0\end{bmatrix},
\alpha_2 = \begin{bmatrix}0 \\ 1 \\ 1\end{bmatrix},
\alpha_3 = \begin{bmatrix}0 \\ 1 \\ -1\end{bmatrix}\] \\
\[Y_1 = \begin{bmatrix}X_1 & X_2 & X_3\end{bmatrix}\alpha_1 = \begin{bmatrix}1 & 1 \\ 0 & -1\end{bmatrix}\]
\[Y_2 = \begin{bmatrix}X_1 & X_2 & X_3\end{bmatrix}\alpha_2 = \begin{bmatrix}0 & 1 \\ 1 & 0\end{bmatrix}\]
\[Y_3 = \begin{bmatrix}X_1 & X_2 & X_3\end{bmatrix}\alpha_3 = \begin{bmatrix}0 & 1 \\ -1 & 0\end{bmatrix}\]
$T$的特征值与线性无关的特征向量为$\lambda_1$、$\lambda_2$、$\lambda_3$，$Y1$、$Y2$、$Y3$

\newpage
\section{第三题}\noindent
计算矩阵多项式 $A^{100}+2A^{50}$，其中
\[
A = \begin{bmatrix}
1 & 1 & -1 \\
1 & 1 & 1 \\
0 & -1 & 2
\end{bmatrix}
\]

\textbf{解：}
\[\mid \lambda E - A \mid= \begin{vmatrix} \lambda - 1 & -1 & 1 \\ -1 & \lambda-1 & -1 \\ 0 & 1 & \lambda-2\end{vmatrix},
\varphi(\lambda) = (\lambda - 1)^2(\lambda - 2)\]
令\[f(\lambda) = \varphi(\lambda)g(\lambda) + r(\lambda)\]
\[\because r(\lambda) < 3 \therefore r(\lambda) = b_0 + b_1\lambda + b_2\lambda^2\]
由$Hamilton-Cayloy$定理 \\
\[\varphi(A)=0 \Rightarrow f(A) = r(A) = b_0 + b_1A + b_2A^2\]
\[\because \varphi(\lambda_i)=0 \therefore f(\lambda_i) = r(\lambda_i)\text{且}f'(\lambda_i)=r'(\lambda_i)\]
带入$\lambda = 2,f(2)=2^{100} + 2 · 2^{50} = r(2) = b_0 + 2b_1 + 4b_2$ \\
带入$\lambda = 1,f(1)=1^{100} + 2 · 1^{50} = r(1) = b_0 + b_1 + b_2$ \\
$f'(\lambda)=100\lambda^99 + 100\lambda^49, f'(1)=200 = b_1 + 2b_2$ \\
\[\Rightarrow \begin{cases}b0 + b1 + b2 = 3 \\ b_1 + 2b_2 = 200 \\ b_0 + 2b_1 + 4b_2 = 2^{100} + 2^{51}\end{cases}
\Rightarrow \begin{cases}b_0 = 2^{100} + 2^{51} - 400 \\ b_2 = -2^{101} - 2^{52} + 606 \\ b_2 = 2^{100} + 2^{51} - 203\end{cases}\]
\[A^ 2 = \begin{bmatrix}2 & 3 & -2 \\ 2 & 1 & 2 \\ -1 & -3 & 3\end{bmatrix}\]
\[f(A) = b_0 + b_1A + b_2A^2\]

\newpage
\section{第四题}\noindent
求矩阵 $A = \begin{bmatrix}
3 & -3 & 2 \\
-1 & 5 & -2 \\
-1 & 3 & 0
\end{bmatrix}$ 的最小多项式。

\textbf{解：}
\[
\begin{aligned}
|\lambda I - A|
&= \begin{vmatrix}
\lambda-3 & 3 & -2 \\
1 & \lambda-5 & 2 \\
1 & -3 & \lambda
\end{vmatrix}
= \begin{vmatrix}
\lambda-3 & 3 & -2 \\
1 & \lambda-5 & 2 \\
0 & 2-\lambda & \lambda-2
\end{vmatrix} \\
&= (\lambda-2) \begin{vmatrix}
\lambda-3 & 3 & -2 \\
1 & \lambda-5 & 2 \\
0 & -1 & 1
\end{vmatrix}
= (\lambda-2) \begin{vmatrix}
\lambda-3 & 1 & -2 \\
1 & \lambda-3 & 2 \\
0 & 0 & 1
\end{vmatrix}
= (\lambda-2) \begin{vmatrix}
\lambda-3 & 1 \\
1 & \lambda-3
\end{vmatrix} \\
&= (\lambda-2)\bigl((\lambda-3)^2-1\bigr) = (\lambda-2)(\lambda^2-6\lambda+8) \\
&= (\lambda-2)^2(\lambda-4)
\end{aligned}
\]
\[\because f_1(A) \ne O\]
\[f_2(A) = (A - 2E)(A-4E)=O\]
所以最小多项式为\[m(\lambda) = (\lambda-2)(\lambda-4)\]

\newpage
\section{第五题}\noindent
设矩阵空间 $\mathbb{R}^{2\times2}$ 的子空间为
\[
V = \{ X = (x_{ij})_{2\times2} \mid x_{11}+x_{12}+x_{21} = 0,\; x_{ij}\in \mathbb{R} \}
\]
$V$ 中的线性变换为
\[
T(X) = X + X^{\mathsf{T}} \quad (\forall X \in V)
\]
求 $V$ 中的一个基，使 $T$ 在该基下的矩阵为对角矩阵。

\textbf{解：}
\[
X = \begin{bmatrix} x_{11} & x_{12} \\ x_{21} & x_{22} \end{bmatrix}
= \begin{bmatrix} -x_{12}-x_{21} & x_{12} \\ x_{21} & x_{22} \end{bmatrix}
= x_{12}\begin{bmatrix} -1 & 1 \\ 0 & 0 \end{bmatrix}
+ x_{21}\begin{bmatrix} -1 & 0 \\ 1 & 0 \end{bmatrix}
+ x_{22}\begin{bmatrix} 0 & 0 \\ 0 & 1 \end{bmatrix}
\]
记
\[
X_1 = \begin{bmatrix} -1 & 1 \\ 0 & 0 \end{bmatrix},\ X_2 = \begin{bmatrix} -1 & 0 \\ 1 & 0 \end{bmatrix},\ X_3 = \begin{bmatrix} 0 & 0 \\ 0 & 1 \end{bmatrix}
\]
$X_1,X_2,X_3$ 线性无关，是 $V$ 的一组基。

\[
\begin{aligned}
T(X_1) &= X_1 + X_1^T = \begin{bmatrix} -1 & 1 \\ 0 & 0 \end{bmatrix} + \begin{bmatrix} -1 & 0 \\ 1 & 0 \end{bmatrix} = \begin{bmatrix} -2 & 1 \\ 1 & 0 \end{bmatrix}
= 1\cdot X_1 + 1\cdot X_2 + 0\cdot X_3, \\[4pt]
T(X_2) &= X_2 + X_2^T = \begin{bmatrix} -1 & 0 \\ 1 & 0 \end{bmatrix} + \begin{bmatrix} -1 & 1 \\ 0 & 0 \end{bmatrix} = \begin{bmatrix} -2 & 1 \\ 1 & 0 \end{bmatrix}
= 1\cdot X_1 + 1\cdot X_2 + 0\cdot X_3, \\[4pt]
T(X_3) &= X_3 + X_3^T = \begin{bmatrix} 0 & 0 \\ 0 & 1 \end{bmatrix} + \begin{bmatrix} 0 & 0 \\ 0 & 1 \end{bmatrix} = \begin{bmatrix} 0 & 0 \\ 0 & 2 \end{bmatrix}
= 0\cdot X_1 + 0\cdot X_2 + 2\cdot X_3.
\end{aligned}
\]
因此 $T$ 在基 $X_1,X_2,X_3$ 下的矩阵满足：
\[
T\begin{bmatrix} X_1 \\ X_2 \\ X_3 \end{bmatrix}
= \begin{bmatrix} 1 & 1 & 0 \\ 1 & 1 & 0 \\ 0 & 0 & 2 \end{bmatrix}
\begin{bmatrix} X_1 \\ X_2 \\ X_3 \end{bmatrix}
\]
记矩阵 $A = \begin{bmatrix} 1 & 1 & 0 \\ 1 & 1 & 0 \\ 0 & 0 & 2 \end{bmatrix}$

\[
|\lambda E - A| = \begin{vmatrix} \lambda - 1 & -1 & 0 \\ -1 & \lambda - 1 & 0 \\ 0 & 0 & \lambda - 2 \end{vmatrix}
= \lambda(\lambda - 2)^2
\]
解得特征值：$\lambda_1 = 0,\ \lambda_{2,3} = 2$，对应的特征向量为：
\[
\alpha_1 = \begin{pmatrix} -1 \\ 1 \\ 0 \end{pmatrix},\ \alpha_2 = \begin{pmatrix} 1 \\ 1 \\ 0 \end{pmatrix},\ \alpha_3 = \begin{pmatrix} 0 \\ 0 \\ 1 \end{pmatrix}
\]
因此对角矩阵与可逆矩阵为：
\[
\Lambda = \begin{bmatrix} 0 & 0 & 0 \\ 0 & 2 & 0 \\ 0 & 0 & 2 \end{bmatrix},\quad
P = \begin{bmatrix} -1 & 1 & 0 \\ 1 & 1 & 0 \\ 0 & 0 & 1 \end{bmatrix}
\]

\vspace{1em}
\noindent (4) 求 $V$ 的基 $Y_1,Y_2,Y_3$：
由 $(Y_1,Y_2,Y_3) = (X_1,X_2,X_3)P$，得：
\[
\begin{aligned}
Y_1 &= -1\cdot X_1 + 1\cdot X_2 + 0\cdot X_3 = \begin{bmatrix} 0 & 1 \\ -1 & 0 \end{bmatrix}, \\[4pt]
Y_2 &= 1\cdot X_1 + 1\cdot X_2 + 0\cdot X_3 = \begin{bmatrix} -2 & 1 \\ 1 & 0 \end{bmatrix}, \\[4pt]
Y_3 &= 0\cdot X_1 + 0\cdot X_2 + 1\cdot X_3 = \begin{bmatrix} 0 & 0 \\ 0 & 1 \end{bmatrix}.
\end{aligned}
\]
故 $T$ 在基 $Y_1,Y_2,Y_3$ 下的矩阵为对角矩阵 $\Lambda$

\newpage
\section{第六题}\noindent
求矩阵 $A$ 的 Jordan 标准形，其中
\[
A = \begin{bmatrix}
-1 & 1 & 0 \\
-4 & 3 & 0 \\
1 & 0 & 2
\end{bmatrix}
\]

\textbf{解：}
\[
|\lambda E - A| = 
\begin{vmatrix}
\lambda + 1 & -1 & 0 \\
4 & \lambda - 3 & 0 \\
-1 & 0 & \lambda - 2
\end{vmatrix}
\]

\[
D_1(\lambda) = 1
\]

因为
\[
\begin{vmatrix}
4 & \lambda - 3 \\
-1 & 0
\end{vmatrix}
= \lambda - 3,
\quad
\begin{vmatrix}
4 & 0 \\
-1 & \lambda - 2
\end{vmatrix}
= 4(\lambda - 2)
\]
$\lambda-3$ 和 $4(\lambda-2)$ 的最大公因式为 $1$，故
\[
D_2(\lambda) = 1
\]

\[
D_3(\lambda) = (\lambda - 2)(\lambda - 1)^2
\]

因此不变因子为：
\[
d_1(\lambda) = 1,\quad d_2(\lambda) = 1,\quad d_3(\lambda) = (\lambda - 2)(\lambda - 1)^2
\]

$\lambda E - A$ 的所有初等因子为：
\[
\lambda - 2,\ (\lambda - 1)^2
\]

故 Jordan 标准形为：
\[
\begin{bmatrix}
2 & 0 & 0 \\
0 & 1 & 1 \\
0 & 0 & 1
\end{bmatrix}
\]


\newpage
\section{第七题}\noindent
已知 $\mathbb{R}^3$ 的线性变换 $T$ 在基 $x_1=(-1,1,1)$，$x_2=(1,0,-1)$，$x_3=(0,1,1)$ 下的矩阵是
\[
\begin{bmatrix}
1 & 0 & 1 \\
1 & 1 & 0 \\
-1 & 2 & 1
\end{bmatrix}
\]
求 $T$ 在基 $e_1=(1,0,0), e_2=(0,1,0), e_3=(0,0,1)$ 下的矩阵。\\
\textbf{解：}
\[
A = \begin{bmatrix}
1 & 1 & 0 \\
1 & 2 & 1 \\
-1 & 1 & 2
\end{bmatrix}
\]

令 \( X = \{x_1, x_2, x_3\} \)，\( E = \{e_1, e_2, e_3\} \)。

由 \( E = XP \)，得 \( P = X^{-1} = \begin{pmatrix}
-1 & 1 & -1 \\
1 & 1 & 0 \\
0 & 0 & 1
\end{pmatrix} \)。

\[
B = XAX^{-1} = \begin{pmatrix}
-1 & 1 & -2 \\
2 & \frac{1}{2} & 0 \\
3 & 0 & 2
\end{pmatrix}
\]

\newpage
\section{第八题}\noindent
给定 $\mathbb{R}^3$ 的两个基
\[
x_1 = (1,0,1),\; x_2 = (2,1,0),\; x_3 = (1,1,1)
\]
\[
y_1 = (1,2,-1),\; y_2 = (2,2,-1),\; y_3 = (2,-1,-1)
\]
定义线性变换 $T x_i = y_i\;(i=1,2,3)$。
\begin{enumerate}
\item 写出由基 $x_1,x_2,x_3$ 到基 $y_1,y_2,y_3$ 的过渡矩阵；
\item 写出 $T$ 在基 $x_1,x_2,x_3$ 下的矩阵；
\item 写出 $T$ 在基 $y_1,y_2,y_3$ 下的矩阵。
\end{enumerate}
\textbf{解：}
(1)
\[
X = \begin{pmatrix}
1 & 2 & 1 \\
0 & 1 & 1 \\
1 & 0 & 1
\end{pmatrix}, \quad
Y = \begin{pmatrix}
1 & 2 & 2 \\
2 & 2 & -1 \\
-1 & -1 & -1
\end{pmatrix}
\]
\[
XP = Y \implies P = X^{-1}Y = \begin{pmatrix}
-2 & -\frac{3}{2} & \frac{3}{2} \\
1 & -\frac{3}{2} & \frac{3}{2} \\
1 & \frac{1}{2} & -\frac{5}{2}
\end{pmatrix}
\]

(2) 令 \( T \) 在基 \( x_1,x_2,x_3 \) 下的矩阵为 \( A \)，则
\[
T(x_1\ x_2\ x_3) = (x_1\ x_2\ x_3)A
\]
由 \( T x_i = y_i \)，得
\[
(x_1\ x_2\ x_3)A = (y_1\ y_2\ y_3) = (x_1\ x_2\ x_3)P
\]
\[
\therefore A = P = \begin{pmatrix}
2 & -\frac{1}{2} & \frac{3}{2} \\
1 & \frac{1}{2} & -\frac{1}{2}
\end{pmatrix}
\]

(3) 令 \( T \) 在基 \( y_1,y_2,y_3 \) 下的矩阵为 \( B \)，则
\[
T(Y) = YB \implies T(XP) = T(X)P = XAP
\]
由 \( XP = Y \implies X = YP^{-1} \)，代入得
\[
T(Y) = YP^{-1}AP = YB
\]
\[
\therefore B = P^{-1}AP
\]
\[
\because A = P, \quad \therefore B = P
\]

\newpage
\section{第九题}\noindent
设 $T$ 是数域 $\mathbb{C}$ 上线性空间 $V^3$ 的线性变换，已知 $T$ 在 $V^3$ 的基 $x_1,x_2,x_3$ 下的矩阵
\[
A = \begin{bmatrix}
3 & 1 & 0 \\
-4 & -1 & 0 \\
4 & -8 & -2
\end{bmatrix}
\]
求 $T$ 的特征值与特征向量。 \\
\textbf{解：}
\[
|\lambda E - A| = 
\begin{vmatrix}
\lambda - 3 & -1 & 0 \\
4 & \lambda + 1 & 0 \\
-4 & 8 & \lambda + 2
\end{vmatrix}
= (\lambda + 2)(\lambda - 1)^2
\]
\[
\therefore \lambda_1 = -2,\ \lambda_{2,3} = 1
\]

\[
\alpha_1 = \begin{pmatrix} 0 \\ 0 \\ 1 \end{pmatrix},\quad 
\alpha_2 = \begin{pmatrix} 3 \\ -6 \\ 20 \end{pmatrix}
\]

\newpage
\section{第十题}\noindent
计算 $2A^8-3A^5+A^4+A^2-4I$，其中
\[
A = \begin{bmatrix}
1 & 0 & 2 \\
0 & -1 & 1 \\
0 & 1 & 0
\end{bmatrix}
\]
\textbf{解：}
\[
\begin{aligned}
f(\lambda) &= |\lambda E - A| = 
\begin{vmatrix}
\lambda - 1 & 0 & 2 \\
0 & \lambda + 1 & -1 \\
0 & -1 & \lambda
\end{vmatrix} 
= (\lambda - 1) \begin{vmatrix}
\lambda + 1 & -1 \\
-1 & \lambda
\end{vmatrix} 
= (\lambda - 1)(\lambda^2 + \lambda - 1) \\
&= \lambda^3 - 2\lambda + 1
\end{aligned}
\]

由 \text{Hamilton–Cayley} 定理，
\[
f(A) = A^3 - 2A + E = O \implies A^3 = 2A - E
\]

\[
\begin{aligned}
&\implies A^4 = 2A^2 - A \\
&A^5 = 2A^3 - A^2 = -A^2 + 4A - 2E \\
&A^8 = 9A^2 - 12A + 4E
\end{aligned}
\]

\[
2A^8 - 3A^5 + A^4 + A^2 - 4E = 24A^2 - 37A + 10E
\]

\[
A^2 = \begin{bmatrix}
1 & 2 & 2 \\
0 & 2 & -1 \\
0 & -1 & 1
\end{bmatrix}
\]

\[
\therefore 2A^8 - 3A^5 + A^4 + A^2 - 4E = 
\begin{bmatrix}
-3 & 48 & -26 \\
0 & 95 & -61 \\
0 & -61 & 34
\end{bmatrix}
\]

\newpage
\section{第十一题}\noindent
求矩阵 $A$ 的最小多项式，其中
\[
A = \begin{bmatrix}
7 & 4 & -4 \\
4 & -8 & -1 \\
-4 & -1 & -8
\end{bmatrix}
\]
\textbf{解：}
\[
\begin{aligned}
|\lambda E - A| &= \begin{vmatrix}
\lambda - 7 & -4 & 4 \\
-4 & \lambda + 8 & 1 \\
4 & 1 & \lambda + 8
\end{vmatrix}
= \begin{vmatrix}
\lambda - 7 & -4 & 4 \\
-4 & \lambda + 8 & 1 \\
0 & \lambda + 9 & \lambda + 9
\end{vmatrix} \\
&= (\lambda + 9) \begin{vmatrix}
\lambda - 7 & -4 & 4 \\
-4 & \lambda + 8 & 1 \\
0 & 1 & 1
\end{vmatrix}
= (\lambda - 9)(\lambda + 9)^2
\end{aligned}
\]

\[
\begin{aligned}
M_{11} &= (\lambda + 7)(\lambda + 9), & M_{12} &= -4(\lambda + 9), & M_{13} &= -4(\lambda + 9) \\
M_{21} &= -4(\lambda + 9), & M_{22} &= (\lambda + 9)(\lambda - 8), & M_{23} &= \lambda + 9 \\
M_{31} &= -4(\lambda + 9), & M_{32} &= \lambda + 9, & M_{33} &= (\lambda + 9)(\lambda - 8)
\end{aligned}
\]

\[
d(\lambda) = \lambda + 9
\]

\[
m(\lambda) = \frac{(\lambda - 9)(\lambda + 9)^2}{\lambda + 9} = (\lambda - 9)(\lambda + 9)
\]

\newpage
\section{第十二题}\noindent
求矩阵 $A$ 的 Jordan 标准形，其中
\[
A = \begin{bmatrix}
3 & 7 & -3 \\
-2 & -5 & 2 \\
-4 & -10 & 3
\end{bmatrix}
\]
\textbf{解：}
\[
\begin{aligned}
|\lambda E - A| &= 
\begin{vmatrix}
\lambda - 3 & -7 & 3 \\
2 & \lambda + 5 & -2 \\
4 & 10 & \lambda - 3
\end{vmatrix}
= 
\begin{vmatrix}
\lambda - 3 & -7 & 3 \\
2 & \lambda + 5 & -2 \\
0 & -2\lambda & \lambda + 1
\end{vmatrix} \\[6pt]
&= 
\begin{vmatrix}
2 & \lambda + 5 & -2 \\
0 & -\lambda^2 - 2\lambda + 1 & 2\lambda \\
0 & -2\lambda & \lambda + 1
\end{vmatrix}
= 2
\begin{vmatrix}
1 & \lambda + 5 & -2 \\
0 & -\lambda^2 - 2\lambda + 1 & 2\lambda \\
0 & -2\lambda & \lambda + 1
\end{vmatrix} \\[6pt]
&= 2
\begin{vmatrix}
-\lambda^2 - 2\lambda + 1 & 2\lambda \\
-2\lambda & \lambda + 1
\end{vmatrix}
= 2(\lambda - 1)(\lambda^2 + 1)
\end{aligned}
\]

\[
\therefore \lambda = 1,\ i,\ -i
\]

\[
\therefore J = \begin{bmatrix}
1 & 0 & 0 \\
0 & i & 0 \\
0 & 0 & -i
\end{bmatrix}
\]

\newpage
\section{第十三题}\noindent
在欧氏空间 $\mathbb{R}^{2\times2}$ 中，矩阵 $A$ 与 $B$ 的内积定义为 $(A,B)=\operatorname{tr}(A^{\mathsf{T}}B)$，子空间
\[
V = \left\{ X = \begin{bmatrix}
x_1 & x_2 \\
x_3 & x_4
\end{bmatrix} \;\middle|\;
\begin{array}{l}
x_1 - x_4 = 0 \\
x_2 - x_3 = 0
\end{array} \right\}
\]
$V$ 中的线性变换为
\[
T(X) = XP + X^{\mathsf{T}} \quad (\forall X \in V),\quad P = \begin{bmatrix}
0 & 1 \\
1 & 0
\end{bmatrix}
\]
\begin{enumerate}
\item 求 $V$ 的一个标准正交基；
\item 验证 $T$ 是 $V$ 中的对称变换；
\item 求 $V$ 的一个标准正交基，使 $T$ 在该基下的矩阵为对角矩阵。
\end{enumerate}
\textbf{解：}
(1)
\[
X = \begin{bmatrix}
x_1 & x_2 \\
x_2 & x_1
\end{bmatrix}
= x_1 \begin{bmatrix}
1 & 0 \\
0 & 1
\end{bmatrix}
+ x_2 \begin{bmatrix}
0 & 1 \\
1 & 0
\end{bmatrix}
= x_1 X_1 + x_2 X_2
\]
取基 \( X_1, X_2 \)，则 \( V \) 的一个标准正交基为
\[
e_1 = \frac{1}{\sqrt{2}} \begin{bmatrix} 1 & 0 \\ 0 & 1 \end{bmatrix}, \quad
e_2 = \frac{1}{\sqrt{2}} \begin{bmatrix} 0 & 1 \\ 1 & 0 \end{bmatrix}
\]

(2). 已知 \( X^T = X \), \( P^T = P \)，则
\[
\langle T(X), Y \rangle
= \operatorname{tr}\left( (XP + X)^T Y \right)
= \operatorname{tr}\left( (PX + X)Y \right)
= \operatorname{tr}\left( PXY + XY \right)
\]
\[
\langle X, T(Y) \rangle
= \operatorname{tr}\left( X^T (YP + Y) \right)
= \operatorname{tr}\left( X(YP + Y) \right)
= \operatorname{tr}\left( XYP + XY \right)
\]

(3).
\[
T(e_1) = \frac{1}{\sqrt{2}} \begin{bmatrix} 1 & 0 \\ 0 & 1 \end{bmatrix} \begin{bmatrix} 0 & 1 \\ 1 & 0 \end{bmatrix} + \begin{bmatrix} 1 & 0 \\ 0 & 1 \end{bmatrix}
= \frac{1}{\sqrt{2}} \begin{bmatrix} 0 & 1 \\ 1 & 0 \end{bmatrix} + \begin{bmatrix} 1 & 0 \\ 0 & 1 \end{bmatrix}
\]
矩阵表示为
\[
\begin{pmatrix} 1 & 1 \\ 1 & 1 \end{pmatrix}, \quad
\lambda_1 = 2,\ \lambda_2 = 0
\]
特征向量
\[
\alpha_1 = \begin{pmatrix} 1 \\ 1 \end{pmatrix}, \quad
\alpha_2 = \begin{pmatrix} 1 \\ -1 \end{pmatrix}
\]
构造标准正交基
\[
f_1 = \frac{1}{\sqrt{2}} (e_1 + e_2), \quad
f_2 = \frac{1}{\sqrt{2}} (e_1 - e_2)
\]
\[
f_1 = \frac{1}{2} \begin{bmatrix} 1 & 1 \\ 1 & 1 \end{bmatrix}, \quad
f_2 = \frac{1}{2} \begin{bmatrix} 1 & -1 \\ -1 & 1 \end{bmatrix}
\]
此为 \( V \) 的一个标准正交基，\( T \) 在该基下的矩阵为
\[
\begin{pmatrix} 2 & 0 \\ 0 & 0 \end{pmatrix}
\]

\end{document}